\documentclass[12pt]{article}
\usepackage{fullpage}
\usepackage{amsmath,amssymb,amsthm}

\newtheorem{lemma}{Lemma}
\newtheorem{proposition}{Proposition}
\newtheorem{assumption}{Assumption}
\newcommand{\R}{\mathbb R}
\newcommand{\E}{\mathbb E}
\newcommand{\1}{\mathbf 1}
\newcommand{\cE}{\mathcal E}
\newcommand{\Var}{\operatorname{Var}}

\begin{document}

\begin{center}
{\bf Invariant measures for the Dirichlet skew stochastic heat equation}
\end{center}

\section*{Problem statement and interpretation}

Let \(E=C_0([0,1])\), and let \(S_t\) be the Dirichlet heat semigroup.
I use the ABLM regularized/mild interpretation specified in the problem,
including the assumed Dirichlet extension and the assumed convergence of
the smooth equations.  The distributional drift has the gradient form
\[
        \delta_0=-\frac12\frac{d}{dr}(-2H_+),\qquad
        H_+(r)=\1_{\{r>0\}} .
\]
Thus the natural reversible measure is a bounded Gibbs perturbation of
Brownian bridge measure.  Let \(\gamma\) be the law on \(E\) of the
standard Brownian bridge from \(0\) to \(0\), and set
\[
   \Lambda(v)=\int_0^1\1_{\{v(x)>0\}}\,dx,\qquad
   \pi(dv)=Z^{-1}e^{2\Lambda(v)}\gamma(dv).       \tag{1}
\]
The map \(\Lambda\) is Borel, since \((v,x)\mapsto v(x)\) is Borel and
Fubini applies.  Also \(0\le \Lambda\le1\), so \(\pi\sim\gamma\), with
density bounded above and below.  The value of \(H_+(0)\) is irrelevant:
\[
   \E_\gamma\int_0^1\1_{\{\beta(x)=0\}}\,dx
    =\int_0^1\mathbb P(\beta(x)=0)\,dx=0 .
\]
The following proves the uniqueness theorem conditional on one
pointwise smoothing input.  The final section explains precisely why
that input is not supplied by the assumptions stated in the problem.

\section*{Consequences of the smooth approximations}

Put
\[
   B_n(r)=\int_{-\infty}^rG_{1/n}(s)\,ds,\qquad
   \pi_n(dv)=Z_n^{-1}\exp\!\left(2\int_0^1B_n(v(x))\,dx\right)\gamma(dv).
\]
Then \(0\le B_n\le1\), \(B_n(r)\to H_+(r)\) for \(r\ne0\), and hence
\[
       \|\pi_n-\pi\|_{\rm TV}\longrightarrow0 .          \tag{2}
\]
Indeed, the normalized densities converge in \(L^1(\gamma)\) by dominated
convergence and by the zero-occupation identity above.

Let \(P_t^n\) be the Markov semigroup of the regularized equation.  For
each fixed \(n\), the Nemytskii drift \(v\mapsto G_{1/n}(v)\) is smooth,
bounded and Lipschitz as an \(L^2(0,1)\)-map.  The standard gradient
SPDE/Dirichlet-form theorem therefore identifies its mild semigroup with
the symmetric diffusion generated, on smooth cylinder functions, by
\[
 L_nF(v)=\frac12{\rm Tr}\,D^2F(v)
       +\left\langle \frac12\Delta_Dv+G_{1/n}(v),DF(v)\right\rangle ,
\]
and by the closed form
\[
   \cE_n(F,G)=\frac12\int_E\langle DF,DG\rangle_{L^2(0,1)}\,d\pi_n .
\]
Equivalently, Gaussian integration by parts gives, first on cylinder
functions and then by closure,
\[
   \int G\,L_nF\,d\pi_n
   =-\frac12\int\langle DF,DG\rangle_{L^2(0,1)}\,d\pi_n
   =\int F\,L_nG\,d\pi_n.                         \tag{3}
\]
Thus \(P_t^n\) is reversible with invariant law \(\pi_n\).

\begin{proposition}\label{prop:limit}
Assume the approximation statement in the problem in the following
semigroup form: for every bounded continuous \(\Phi:E\to\R\), every
\(x\in E\), and every \(t>0\),
\[
        P_t^n\Phi(x)\longrightarrow P_t\Phi(x).      \tag{4}
\]
Then \(\pi\) is invariant for \(P_t\), \(P_t\) is symmetric on
\(L^2(\pi)\), and the \(L^2(\pi)\)-semigroup has a spectral gap.
\end{proposition}

\begin{proof}
For \(f,g\in C_b(E)\), dominated convergence with respect to \(\pi\),
(2), and invariance of \(\pi_n\) give
\[
 \int P_tf\,d\pi
 =\lim_n\int P_t^nf\,d\pi
  =\lim_n\int P_t^nf\,d\pi_n
  =\lim_n\int f\,d\pi_n
  =\int f\,d\pi .
\]
Similarly, using (3),
\[
 \int g\,P_tf\,d\pi
 =\lim_n\int g\,P_t^nf\,d\pi_n
  =\lim_n\int f\,P_t^ng\,d\pi_n
  =\int f\,P_tg\,d\pi .
\]
Since \(C_b(E)\) is dense in \(L^2(\pi)\), and invariance makes \(P_t\) an
\(L^2(\pi)\)-contraction, this extends to all of \(L^2(\pi)\).

The Brownian bridge satisfies the Gaussian Poincare inequality
\[
     \Var_\gamma(F)\le C_0\int\|DF\|_{L^2(0,1)}^2\,d\gamma
\]
on cylinder functions.  Because the densities \(d\pi_n/d\gamma\) are
bounded above and below uniformly in \(n\), there is \(C<\infty\) such
that
\[
     \Var_{\pi_n}(F)\le C\,\cE_n(F,F) .             \tag{5}
\]
By the spectral theorem for the reversible semigroup generated by
\(\cE_n\), (5) implies, with \(c>0\) independent of \(n\),
\[
     \Var_{\pi_n}(P_t^n f)
       \le e^{-2ct}\Var_{\pi_n}(f),\qquad f\in L^2(\pi_n).   \tag{6}
\]
For bounded continuous \(f\), (4) and (2) allow passage to the limit in
(6), yielding
\[
     \Var_{\pi}(P_t f)\le e^{-2ct}\Var_\pi(f).       \tag{7}
\]
Density and \(L^2(\pi)\)-contractivity extend (7) to every
\(f\in L^2(\pi)\).
\end{proof}

\section*{Uniqueness once pointwise smoothing is known}

\begin{lemma}\label{lem:criterion}
Assume Proposition \ref{prop:limit}.  If, for some \(\lambda>0\),
\[
     R_\lambda^P(x,\cdot)\ll \pi,\qquad
     R_\lambda^P(x,A):=\int_0^\infty e^{-\lambda t}P_t(x,A)\,dt,
     \quad x\in E,                                      \tag{8}
\]
then \(\pi\) is the only invariant probability measure of \(P_t\).
\end{lemma}

\begin{proof}
Let \(\eta\) be invariant.  If \(\pi(A)=0\), then (8) and invariance give
\[
   \eta(A)=\lambda\int_E R_\lambda^P(x,A)\,\eta(dx)=0.
\]
Thus \(\eta=h\pi\).  For bounded \(\varphi\),
\[
 \int \varphi h\,d\pi=\int P_t\varphi\,h\,d\pi
   =\int \varphi\,P_t h\,d\pi ,
\]
where the last equality follows first for bounded \(h\) by symmetry and
then for \(h\in L^1(\pi)\) by truncation.  Hence \(P_t h=h\) in
\(L^1(\pi)\).

Fix \(a>0\), and set \(h_a=h\wedge a\).  Since \(r\mapsto r\wedge a\) is
concave and \(P_t\) is Markov,
\[
      h_a=(P_t h)\wedge a\ge P_t(h\wedge a)=P_t h_a .
\]
The two sides have the same \(\pi\)-integral, so equality holds
\(\pi\)-a.s.  Now \(h_a\in L^2(\pi)\), and (7) gives
\[
      \Var_\pi(h_a)=\Var_\pi(P_t h_a)
          \le e^{-2ct}\Var_\pi(h_a).
\]
For any \(t>0\), this forces \(\Var_\pi(h_a)=0\).  Thus \(h\wedge a\) is
constant for every \(a>0\).  Varying \(a\) shows that \(h\) itself is
constant; since \(\int h\,d\pi=1\), \(h=1\) and \(\eta=\pi\).
\end{proof}

Here are two useful sufficient forms of the missing input.

\begin{lemma}\label{lem:sf}
If \(P_t\) is strong Feller for every \(t>0\), then (8) holds.
The same conclusion holds if, for one \(\lambda>0\), the resolvent
\(R_\lambda^P\) maps bounded Borel functions to continuous functions.
\end{lemma}

\begin{proof}
Suppose first that \(P_t\) is strong Feller.  Let \(A\) be Borel with
\(\pi(A)=0\).  For fixed \(t>0\), the function
\(\varphi_t(x)=P_t(x,A)\) is continuous and nonnegative.  Invariance of
\(\pi\) gives \(\int\varphi_t\,d\pi=\pi(A)=0\).  Since \(\pi\sim\gamma\)
and Brownian bridge measure has full support on \(E\), a nonnegative
continuous function vanishing \(\pi\)-a.e. is identically zero.  Hence
\(P_t(x,A)=0\) for all \(x\) and all \(t>0\), and (8) follows by
integration in \(t\).  The resolvent-strong-Feller version is identical:
\(x\mapsto R_\lambda^P(x,A)\) is continuous, has zero \(\pi\)-integral,
and therefore is identically zero.
\end{proof}

\section*{What Bounebache--Zambotti gives}

For the bounded-variation potential \(f=-2H_+\), the
Bounebache--Zambotti Dirichlet form is exactly
\[
        \cE(F,G)=\frac12\int_E\langle DF,DG\rangle_{L^2(0,1)}\,d\pi .
\]
Their Proposition 2.5 gives an \(L^2(\pi)\)-resolvent kernel
\(\rho_\lambda(x,dy)\) with \(\rho_\lambda(x,\cdot)\ll\pi\) for all
\(x\).  However the identity of this kernel with the Markov-process
resolvent is asserted only \(\pi\)-a.e.; in the construction it is
extended over an exceptional set by definition.  Their local-time SPDE is
therefore obtained for quasi-every initial condition.  Combining that
with ABLM weak/mild uniqueness identifies the two constructions
quasi-everywhere, but not automatically at every deterministic
\(x\in E\).

One exact way to bridge this is the following entrance statement.

\begin{lemma}\label{lem:entrance}
Let \(N\) be a Bounebache--Zambotti properly exceptional set such that,
for every \(y\notin N\), \(R_\lambda^P(y,\cdot)\ll\pi\).  If
\[
          P_s(x,N)=0,\qquad x\in E,\ s>0,              \tag{9}
\]
then (8) holds for every \(x\in E\).
\end{lemma}

\begin{proof}
Let \(\pi(A)=0\).  For \(y\notin N\), \(R_\lambda^P(y,A)=0\).  For any
\(x\in E\) and \(\varepsilon>0\), the Markov property gives
\[
 R_\lambda^P(x,A)
 =\int_0^\varepsilon e^{-\lambda t}P_t(x,A)\,dt
  +e^{-\lambda\varepsilon}\int_E R_\lambda^P(y,A)P_\varepsilon(x,dy).
\]
The first term is at most \(\varepsilon\).  The second term is bounded by
\(\lambda^{-1}P_\varepsilon(x,N)\), hence is zero by (9).  Letting
\(\varepsilon\downarrow0\) proves \(R_\lambda^P(x,A)=0\).
\end{proof}

\section*{Remaining open issues}

The proof above shows that the desired statement is true once the actual
ABLM Dirichlet semigroup has any one of the pointwise smoothing
properties in Lemma \ref{lem:criterion}, Lemma \ref{lem:sf}, or Lemma
\ref{lem:entrance}.  What remains unproved from the hypotheses stated in
the problem is precisely
\[
\boxed{\quad
   R_\lambda^P(x,\cdot)\ll\pi
   \quad\hbox{for every deterministic }x\in C_0([0,1]).
   \quad}
\]
Mathematically, this is needed only to exclude invariant probabilities
carried by a \(\pi\)-null exceptional set.  The smooth approximation
argument gives the invariant reversible measure \(\pi\) and an
\(L^2(\pi)\) spectral gap, while Bounebache--Zambotti gives the required
absolute continuity only for the \(L^2(\pi)\) version, equivalently away
from an exceptional set.  To close the proof one would need, for example:

\begin{enumerate}
\item a strong-Feller or resolvent-strong-Feller theorem for the
Dirichlet ABLM semigroup;
\item a uniform \(L^p\), heat-kernel, or capacity estimate for the
regularized semigroups \(P_t^n\) that survives the limit \(n\to\infty\);
\item a proof that the ABLM process immediately enters the
Bounebache--Zambotti non-exceptional set, i.e. (9);
\item or a strict all-starting-points identification of the ABLM process
with a version of the Bounebache--Zambotti Dirichlet-form process whose
resolvent kernel represents the true resolvent at every \(x\).
\end{enumerate}

Without one of these additional inputs, the argument proves uniqueness
only among invariant measures absolutely continuous with respect to
\(\pi\), and proves full uniqueness conditionally on the boxed statement.

\begin{thebibliography}{9}
\bibitem{ABLM}
S. Athreya, O. Butkovsky, K. L\^e and L. Mytnik,
\emph{Well-posedness of stochastic heat equation with distributional drift
and skew stochastic heat equation}, Comm. Pure Appl. Math. 77 (2024),
2708--2777.

\bibitem{ABLMweakmild}
S. Athreya, O. Butkovsky, K. L\^e and L. Mytnik,
\emph{Analytically weak and mild solutions to stochastic heat equation with
irregular drift}, Stoch. PDE: Anal. Comp. (2025).

\bibitem{BZ}
S. K. Bounebache and L. Zambotti,
\emph{A skew stochastic heat equation}, J. Theoret. Probab. 27 (2014),
168--201.
\end{thebibliography}

\end{document}
